\documentclass[11pt,a4paper]{article}

\usepackage[utf8]{inputenc}
\usepackage[T2A]{fontenc}
\usepackage[ukrainian,english]{babel}
\usepackage[top=22mm,bottom=22mm,left=22mm,right=22mm]{geometry}
\usepackage{hyperref}
\usepackage{fancyhdr}
\usepackage{parskip}
\usepackage{tcolorbox}
\usepackage{enumitem}
\usepackage{listings}
\usepackage{xcolor}

\tcbuselibrary{skins,breakable}

\hypersetup{
    pdftitle={Workshop 12 --- Debugging Claude Code: Handout},
    pdfauthor={Vadym Bondarenko},
    colorlinks=true,
    linkcolor=blue,
    urlcolor=cyan,
}

\pagestyle{fancy}
\fancyhf{}
\fancyhead[L]{\small\textit{Workshop 12 --- Debugging}}
\fancyhead[R]{\small\thepage}
\fancyfoot[C]{\small\textit{vcdev.me \textperiodcentered{} 2026}}
\renewcommand{\headrulewidth}{0.4pt}

\lstdefinelanguage{json}{
    keywords={true,false,null},
    sensitive=false,
    morestring=[b]",
    morecomment=[l]{//}
}

\lstdefinelanguage{yaml}{
    keywords={true,false,null},
    sensitive=false,
    comment=[l]{\#},
    morestring=[b]',
    morestring=[b]"
}

\lstset{
    basicstyle=\ttfamily\small,
    breaklines=true,
    frame=single,
    backgroundcolor=\color{gray!10},
    rulecolor=\color{gray!40},
    xleftmargin=8pt,
    xrightmargin=8pt,
    aboveskip=6pt,
    belowskip=6pt,
}

\title{%
  {\Huge\bfseries Workshop 12: Debugging Claude Code}\\[0.6em]
  {\large Коли ламається --- handout}
}
\author{Vadym Bondarenko \\ \small\href{mailto:vadym.bondarenko@vcdev.me}{vadym.bondarenko@vcdev.me}}
\date{2026}

\begin{document}

\maketitle

\begin{abstract}
Пост-воркшоп референс. Не повторює слайди --- сухе зведення команд, схем, чек-листів і jq-запитів. Усі факти посилаються на офіційні docs (\href{https://code.claude.com/docs/en/troubleshooting}{code.claude.com/docs/en/troubleshooting}, \href{https://code.claude.com/docs/en/debug-your-config}{debug-your-config}, \href{https://code.claude.com/docs/en/hooks}{hooks}, \href{https://code.claude.com/docs/en/skills}{skills}, \href{https://code.claude.com/docs/en/mcp}{mcp}).
\end{abstract}

\tableofcontents
\newpage

\section{Diagnostic surface}

8 команд, які показують, що насправді завантажилось у поточну сесію.

\begin{tabular}{ll}
\hline
\textbf{Команда} & \textbf{Що показує} \\
\hline
\texttt{/context} & Усе у контекстному вікні (system, memory, skills, MCP) \\
\texttt{/doctor} & Конфіг-валідація: invalid keys, schema errors \\
\texttt{/status} & Активні settings sources, auth method \\
\texttt{/memory} & Завантажені CLAUDE.md / rules \\
\texttt{/skills} & Available skills + source \\
\texttt{/hooks} & Зареєстровані hooks за event \\
\texttt{/mcp} & MCP-сервери: connection status, tools \\
\texttt{/permissions} & Resolved allow/deny rules \\
\hline
\end{tabular}

\textbf{Порядок дій:} \texttt{/context} $\to$ конкретна команда $\to$ \texttt{--debug <category>} якщо неясно.

\section{Hook debugging}

\subsection{Exit codes}

\begin{tabular}{lll}
\hline
\textbf{Code} & \textbf{Поведінка} & \textbf{JSON оброблюється?} \\
\hline
0 & Success --- дія дозволена & Так (stdout як decision) \\
2 & \textbf{Blocking} --- дія заблокована, stderr $\to$ Claude & Ні \\
Інше & Non-blocking --- дія проходить, перший рядок stderr $\to$ юзеру & Ні \\
\hline
\end{tabular}

\textbf{Pitfall:} \texttt{exit 1} НЕ блокує. Це non-blocking error. Для блокування --- \texttt{exit 2}.

\subsection{Exit 2 ефекти за event}

\begin{lstlisting}
PreToolUse              -> блокує tool call
UserPromptSubmit        -> блокує prompt, стирає з контексту
UserPromptExpansion     -> блокує expansion
PermissionRequest       -> denies permission
Stop                    -> не дає Claude зупинитись
WorktreeCreate          -> non-zero валить creation
\end{lstlisting}

PostToolUse, SessionEnd --- exit 2 лише показує stderr юзеру, заблокувати не може (дія вже сталась).

\subsection{Stderr/stdout}

\begin{tabular}{ll}
\hline
\textbf{Сценарій} & \textbf{Куди йде} \\
\hline
Exit 0, plain text stdout & Debug log only (НЕ transcript) \\
Exit 0, JSON stdout & Парситься як decision \\
Exit 2, stderr & \textbf{Claude бачить} як error \\
Exit 1+, stderr & Перший рядок $\to$ transcript, повний $\to$ debug log \\
\hline
\end{tabular}

3 події що ВКИДАЮТЬ stdout у контекст: \texttt{SessionStart}, \texttt{UserPromptSubmit}, \texttt{UserPromptExpansion}. Cap: \textbf{10,000 символів}.

\subsection{JSON output schema}

Universal fields (для всіх events):

\begin{lstlisting}[language=json]
{
  "continue": true,
  "stopReason": "Build failed",
  "suppressOutput": false,
  "systemMessage": "Warning to user"
}
\end{lstlisting}

PreToolUse hookSpecificOutput:

\begin{lstlisting}[language=json]
{
  "hookSpecificOutput": {
    "hookEventName": "PreToolUse",
    "permissionDecision": "allow|deny|ask|defer",
    "permissionDecisionReason": "reason text",
    "updatedInput": { "command": "..." },
    "additionalContext": "context for Claude"
  }
}
\end{lstlisting}

\subsection{\texttt{/hooks} і \texttt{--debug hooks}}

\textbf{\texttt{/hooks}} --- список усіх зареєстрованих hooks за event. Якщо твого нема:

\begin{itemize}
\item \texttt{matcher} --- JSON array замість string. Має бути \texttt{"Edit\textbar{}Write"}, не \texttt{["Edit", "Write"]}
\item \texttt{matcher} --- lowercase. Case-sensitive: \texttt{Bash}, \texttt{Edit}, \texttt{Write}, \texttt{Read}
\item Конфіг у \texttt{.claude/hooks.json} (нема такого файлу) $\to$ перенеси у \texttt{settings.json}, ключ \texttt{"hooks"}
\item Schema error $\to$ \texttt{/doctor} репортує, entry дропається тихо
\end{itemize}

\textbf{\texttt{claude --debug hooks}} --- live-трасування подій:

\begin{itemize}
\item Кожен event evaluated
\item Які matchers перевірились
\item Hook exit code
\item Повний output (не лише перший рядок stderr)
\end{itemize}

\subsection{Pitfall: shell profile pollution}

stdout мусить мати ТІЛЬКИ JSON. Якщо \texttt{\textasciitilde/.bashrc} echo-ить welcome --- JSON не запарситься.

\begin{lstlisting}[language=bash]
#!/bin/bash
exec 2>/dev/null
source ~/.profile
exec 2>&1
jq -n '{ "continue": true }'
\end{lstlisting}

\section{Skill debugging --- 6-крокова перевірка}

\begin{enumerate}
\item \textbf{\texttt{/}-меню видно skill?} Ні $\to$ перевір шлях \texttt{\textasciitilde/.claude/skills/<name>/SKILL.md}.
\item \textbf{Перефразуй точно під description.} Спрацювало $\to$ tuning description.
\item \textbf{Грепни \texttt{disable-model-invocation} / \texttt{user-invocable}}.
\item \textbf{Перевір \texttt{paths}} --- може skill активується лише на матчинг файлах.
\item \textbf{Manual \texttt{/skill-name} працює?} Так $\to$ skill loaded, проблема matching.
\item \textbf{Запит надто простий?} Claude не тригерить skill для тривіальних задач (\textit{by design}).
\end{enumerate}

\subsection{Description budget}

$\sim$1\% контекстного вікна $=$ $\sim$8K символів за замовчуванням. Cap per entry: \textbf{1,536 символів}. Багато skill-ів $\to$ описи обрізаються.

\textbf{Підняти:} \texttt{SLASH\_COMMAND\_TOOL\_CHAR\_BUDGET=16000}

\subsection{Lifecycle після compaction}

\begin{itemize}
\item Most recent invocation кожного skill-а --- preserved
\item Перших 5,000 токенів кожного збережено
\item Сумарний бюджет: \textbf{25,000 токенів}
\item Старіші invocation-и можуть дропнутись повністю
\end{itemize}

\textbf{Mitigation:} re-invoke skill якщо «забувся».

\section{MCP debugging}

\subsection{\texttt{/mcp} статуси}

\begin{tabular}{ll}
\hline
\textbf{Статус} & \textbf{Що означає} \\
\hline
\texttt{connected} & Сервер працює, tools доступні \\
\texttt{pending approval} & \texttt{.mcp.json} сервер чекає approve \\
\texttt{failed} & Стартував і помер. Зазвичай: relative path \\
\texttt{reconnecting} & HTTP/SSE розрив, exp backoff (5 спроб) \\
\texttt{connected (0 tools)} & Стартував, але tools/list порожній \\
\hline
\end{tabular}

\subsection{\texttt{.mcp.json} schema}

Лежить у \textbf{корені проєкту} (НЕ у \texttt{.claude/}).

\begin{lstlisting}[language=json]
{
  "mcpServers": {
    "github": {
      "command": "npx",
      "args": ["-y", "@modelcontextprotocol/server-github"],
      "env": { "GITHUB_TOKEN": "${GH_TOKEN}" }
    },
    "api": {
      "type": "http",
      "url": "${API_BASE_URL:-https://api.example.com}/mcp",
      "headers": {
        "Authorization": "Bearer ${API_KEY}"
      }
    }
  }
}
\end{lstlisting}

\textbf{Env expansion:} \texttt{\$\{VAR\}}, \texttt{\$\{VAR:-default\}}. Required без default $\to$ fail to parse.

\subsection{Топ-5 фейлів}

\begin{tabular}{p{5cm}p{4cm}p{5cm}}
\hline
\textbf{Симптом} & \textbf{Причина} & \textbf{Фікс} \\
\hline
\texttt{.mcp.json} ігнорується & Файл у \texttt{.claude/} & Корінь проєкту \\
Сервер не з'являється & Approval prompt dismissed & \texttt{/mcp} $\to$ Approve \\
Fail у деяких теках & Relative path & Абсолютний шлях \\
Env vars missing & Vars у \texttt{settings.json} & Per-server \texttt{env} у \texttt{.mcp.json} \\
\texttt{Failed to parse config} & Required \texttt{\$\{VAR\}} not set & Set env або \texttt{\$\{VAR:-default\}} \\
\hline
\end{tabular}

\subsection{\texttt{claude --debug mcp}}

Логує: handshake (\texttt{initialize}, \texttt{initialized}), сервер stderr, \texttt{tools/list} response, reconnect attempts.

\subsection{Транспорти}

\begin{itemize}
\item \textbf{stdio} --- локальний процес. Не auto-reconnect.
\item \textbf{HTTP} --- recommended для remote. Auto-reconnect: 5 спроб, 1s$\to$2s$\to$4s$\to$8s$\to$16s.
\item \textbf{SSE} --- DEPRECATED. Use HTTP instead.
\end{itemize}

\subsection{OAuth pitfalls}

\begin{itemize}
\item \emph{Incompatible auth server: does not support dynamic client registration} $\to$ зареєструй OAuth-app вручну, подай creds у \texttt{claude mcp add}
\item Browser redirect failed $\to$ скопіюй callback URL, встав у CC prompt
\item Internal SSO/Kerberos/short-lived $\to$ \texttt{headersHelper} у \texttt{.mcp.json}
\end{itemize}

\section{Session transcripts}

\subsection{Локація}

\begin{lstlisting}
~/.claude/projects/<project-slug>/<session-id>.jsonl
\end{lstlisting}

\texttt{<project-slug>} --- sanitized cwd path (наприклад \texttt{-home-vadym-projects-foo}).

Кожен рядок --- JSON event сесії: user prompts, assistant responses, tool calls/results, hook outputs, system messages, compact boundaries.

\textbf{Caveat:} формальна schema не задокументована. Поля (\texttt{.type}, \texttt{.message.content}, \texttt{.subtype}) спостережено в реальних transcripts але не у specs.

\subsection{jq queries --- основні}

\begin{lstlisting}[language=bash]
# Скільки events
wc -l session.jsonl

# Які типи зустрічаються
jq -r '.type' session.jsonl | sort -u

# Усі user prompts
jq 'select(.type == "user") | .message.content' session.jsonl

# Усі tool calls
jq 'select(.type == "assistant") | .message.content[]?
    | select(.type == "tool_use")' session.jsonl

# Compaction boundaries
jq 'select(.type == "system" and .subtype == "compact_boundary")' session.jsonl

# Hook outputs
jq 'select(.type == "system" and ((.subtype // "") | startswith("hook_")))' session.jsonl

# Останні N events
tail -10 session.jsonl | jq .
\end{lstlisting}

\section{Auto-compaction}

\subsection{Preserved}

\begin{itemize}
\item Сесійний summary (Claude-generated)
\item Останні N turn-ів verbatim
\item System prompt + memory --- re-injected fresh
\item Skills --- most recent each, first 5K tokens, $\Sigma$ 25K budget
\end{itemize}

\subsection{Lost}

\begin{itemize}
\item Старі tool outputs verbatim
\item Earlier file reads (тільки summary references)
\item Skills не invoked недавно АБО витіснені 25K-бюджетом
\item Детальні code chunks середини сесії
\end{itemize}

\subsection{Mitigation --- 5 тактик}

\begin{enumerate}
\item Re-invoke skill (контент завантажиться знову)
\item \texttt{/recap} (silent injection) перед заповненням контексту
\item Subagents (\texttt{context: fork} у skill)
\item \texttt{/compact <focus>} з directive (наприклад «keep only the plan»)
\item \texttt{/clear} якщо стара частина не потрібна
\end{enumerate}

\subsection{Autocompact thrashing}

Симптом: \texttt{Autocompact is thrashing: the context refilled to the limit}. Файл/tool output постійно перезаповнює контекст.

\textbf{Фікс:} читай файл chunks, drop large output, або \texttt{/clear}.

\section{Common errors}

\begin{tabular}{p{5cm}p{4.5cm}p{4.5cm}}
\hline
\textbf{Error} & \textbf{Причина} & \textbf{Фікс} \\
\hline
Hook exited code 1 (не блокує) & exit 1 замість 2 & \texttt{exit 2} \\
Hook JSON parse error & non-JSON у stdout & \texttt{exec 2>/dev/null} \\
Skill not found & flat \texttt{.md} замість теки & \texttt{<name>/SKILL.md} \\
MCP failed to start & relative path / no exec & абсолютний шлях; \texttt{--debug mcp} \\
MCP: Connection refused & endpoint мертвий / stdio process помер & test endpoint \\
Autocompact is thrashing & файл reflood-ить контекст & chunks; \texttt{/compact}; \texttt{/clear} \\
\texttt{spawn claude ENOENT} (CC як MCP) & Wrong path до \texttt{claude} & абсолютний шлях у \texttt{command} \\
\hline
\end{tabular}

\section{Production debug-checklist}

\begin{itemize}[leftmargin=2em]
\item[$\square$] \texttt{/doctor} --- перевір спочатку
\item[$\square$] \texttt{/context} --- що завантажилось у window
\item[$\square$] \texttt{/hooks} / \texttt{/skills} / \texttt{/mcp} --- конкретний компонент
\item[$\square$] Перевір scope precedence: managed $>$ user $>$ project $>$ local
\item[$\square$] \texttt{--debug <category>} у нову сесію
\item[$\square$] Transcript: \texttt{\textasciitilde/.claude/projects/<slug>/*.jsonl} через jq
\item[$\square$] Re-invoke skill якщо post-compaction
\item[$\square$] Restart CC для нових тек skill / новий plugin
\end{itemize}

\section{Команди для вправ}

\subsection{Вправа 1 --- broken hook}

\begin{lstlisting}[language=bash]
mkdir -p /tmp/hook-debug && cd /tmp/hook-debug
mkdir -p .claude
cp <starter>/settings.json .claude/settings.json
cp <starter>/block-rm.sh .claude/block-rm.sh
chmod +x .claude/block-rm.sh
claude --debug hooks
\end{lstlisting}

\subsection{Вправа 3 --- broken \texttt{.mcp.json}}

\begin{lstlisting}[language=bash]
mkdir -p /tmp/mcp-debug && cd /tmp/mcp-debug
cp <starter>/.mcp.json .
cp <starter>/echo-server.py .
chmod +x echo-server.py
claude --debug mcp
\end{lstlisting}

\subsection{Вправа 4 --- transcript jq}

\begin{lstlisting}[language=bash]
cd 04-transcript-jq
bash ../solutions/04-transcript-jq/queries.sh starter/session.jsonl
\end{lstlisting}

\section{Resources}

\begin{itemize}
\item \href{https://code.claude.com/docs/en/troubleshooting}{code.claude.com/docs/en/troubleshooting} --- install/auth/perf/IDE
\item \href{https://code.claude.com/docs/en/debug-your-config}{code.claude.com/docs/en/debug-your-config} --- \texttt{/context}, \texttt{/doctor}, common-cause table
\item \href{https://code.claude.com/docs/en/hooks}{code.claude.com/docs/en/hooks} --- exit codes, JSON schema, \texttt{--debug hooks}
\item \href{https://code.claude.com/docs/en/skills}{code.claude.com/docs/en/skills} --- troubleshooting, lifecycle, budget
\item \href{https://code.claude.com/docs/en/mcp}{code.claude.com/docs/en/mcp} --- \texttt{.mcp.json}, transports, OAuth
\item Цей репо: \texttt{workshops/12-debugging/exercises/} і \texttt{solutions/}
\end{itemize}

\end{document}
